% demo_showcase.tex — LaTeX Optimizer Skill 演示文档
% 展示大白话 → 专业排版的完整效果

\documentclass[12pt,a4paper]{ctexart}

% === 导言区：由 Skill 自动优化 ===
\usepackage{amsmath,amssymb}
\usepackage{booktabs}
\usepackage{geometry}
\usepackage{setspace}
\usepackage{xcolor}
\geometry{margin=2.5cm}
\onehalfspacing

\title{\bfseries EasyTex Agent 功能演示}
\author{LaTeX Optimizer Skill}
\date{\today}

\begin{document}
\maketitle
\tableofcontents
\newpage

% ============================================
\section{大白话转公式：矩阵乘法}
% ============================================

\textbf{用户输入：} ``写一个3x3单位矩阵右乘列向量x，等号右边是全是lambda的向量''

\textbf{Skill 生成的 LaTeX 代码：}

\begin{equation}
    I_3 \cdot \mathbf{x} =
    \begin{bmatrix}
        1 & 0 & 0 \\
        0 & 1 & 0 \\
        0 & 0 & 1
    \end{bmatrix}
    \cdot
    \begin{bmatrix}
        x_1 \\ x_2 \\ x_3
    \end{bmatrix}
    =
    \begin{bmatrix}
        \lambda \\ \lambda \\ \lambda
    \end{bmatrix}
    = \lambda \cdot
    \begin{bmatrix}
        1 \\ 1 \\ 1
    \end{bmatrix}
    \label{eq:identity}
\end{equation}

公式~\eqref{eq:identity} 展示了单位矩阵 $I_3$ 作用于任意向量 $\mathbf{x}$ 时，
结果等价于将 $\mathbf{x}$ 投影到全一向量方向并缩放 $\lambda$ 倍。

% ============================================
\section{大白话转公式：协方差矩阵}
% ============================================

\textbf{用户输入：} ``3x3的协方差矩阵Sigma，对角线是sigma_1^2, sigma_2^2, sigma_3^2，非对角线是rho_ij乘对应sigma''

\textbf{Skill 输出：}

\begin{equation}
    \boldsymbol{\Sigma} =
    \begin{bmatrix}
        \sigma_1^2 & \rho_{12}\sigma_1\sigma_2 & \rho_{13}\sigma_1\sigma_3 \\[4pt]
        \rho_{21}\sigma_2\sigma_1 & \sigma_2^2 & \rho_{23}\sigma_2\sigma_3 \\[4pt]
        \rho_{31}\sigma_3\sigma_1 & \rho_{32}\sigma_3\sigma_2 & \sigma_3^2
    \end{bmatrix}
    \label{eq:cov}
\end{equation}

% ============================================
\section{大白话转公式：贝叶斯定理}
% ============================================

\textbf{用户输入：} ``写贝叶斯定理，P(A|B)等于P(B|A)P(A)除以P(B)''

\textbf{Skill 输出：}

\begin{equation}
    P(A \mid B) = \frac{P(B \mid A) \cdot P(A)}{P(B)}
    \label{eq:bayes}
\end{equation}

其中 $P(B) = \sum_{i} P(B \mid A_i) P(A_i)$ 为全概率公式展开。

% ============================================
\section{排版优化展示}
% ============================================

\textbf{优化前}（默认 Computer Modern，单倍行距）$\longrightarrow$
\textbf{优化后}（中文思源宋体 / 英文 Times New Roman，1.5 倍行距）。

\begin{table}[h]
\centering
\caption{Skill 自动修改导言区对照表}
\begin{tabular}{@{}lll@{}}
\toprule
\textbf{修改项} & \textbf{修改前} & \textbf{修改后} \\
\midrule
正文字体   & Computer Modern    & Times New Roman \\
中文字体   & 未设置             & Source Han Serif SC \\
行间距     & 单倍 (1.0)        & 1.5 倍 \\
页边距     & 默认 (3.2cm)      & 2.5cm \\
\bottomrule
\end{tabular}
\end{table}

% ============================================
\section{编译验证闭环}
% ============================================

当用户通过自然语言指令修改文档后，Skill 自动执行：

\begin{enumerate}
    \item \textbf{解析导言区} — 识别已有宏包，避免重复加载
    \item \textbf{字体预检查} — 确保指定字体已在系统中安装
    \item \textbf{安全修改} — 备份原文件，注入新配置
    \item \textbf{静默编译} — 后台运行 \texttt{xelatex}，捕获全部输出
    \item \textbf{错误过滤} — 从 .log 中提取精准报错，过滤警告噪声
\end{enumerate}

\textbf{示例：} 若用户指定了系统中不存在的字体，Skill 返回：

\begin{verbatim}
FATAL → FONT_NOT_FOUND
→ 停止修改，提示用户安装字体或选择替代字体
→ 原文件未被修改
\end{verbatim}

\end{document}
